\begin{abstract}
The European Union's Ecodesign for Sustainable Products Regulation (\ESPR) mandates Digital Product Passports (\DPPs) for an expanding range of product categories. The current architectural blueprint envisions flat-file data payloads hosted at the Economic Operator level, accessed via unique identifiers resolving to distributed endpoints. This paper argues that this model carries systemic vulnerabilities --- link rot, business continuity failure, cyberattack exposure, and intellectual property leakage --- that will compromise the \ESPR's circular economy objectives over the decades-long lifespan envisioned for \DPP\ data. We propose a portal-backed, graph-based,
network-structured \DPP\ architecture built on semantic web standards (\RDF, \JSONLD, Named Graphs), in which each individual \DPP\ is a knowledge graph node in a production network, linked to its predecessors via semantic properties. All product data --- including data from the Responsible Economic Operator --- is stored behind a federated portal, with strictly controlled, branch-level access rights ensuring
protection of sensitive data. \SHACL\ serves as a unified data modelling language for both conformity checking and user-interface generation, supported by a public repository of ontologised regulations and a modular \SHACL\ brick library. Evidence from the \DPPforEU~2026 conference confirms that the constituent elements of this architecture are independently recognised by the practitioner community as necessary. We propose to integrate them into a coherent whole.
\end{abstract}
